\scanpage{185}
De plus, on a une structure \unsure{polarisée|polaris.|privilégiée}\inserted{ $\pi$} sur $G_{\mathbf R}^{\circ}$, qu'on peut \deleted{\unsure{défi|définir}} supposer définie par un hom \deleted{$\unsure{\mu_2\to G}$} $i:\mu_{2\,\mathbf C}\longrightarrow G_{\mathbf C}^{\circ}$ (comme on voit en envoyant un tore \inserted{\unsure{provenant d'un $s$|prenant un sens}} dedans). D'autre part, $G(s')_{\mathbf C}$ est muni de $i(s'):G_{m\,\mathbf C}\to G(s')_{\mathbf C}$, et \unsure{on|qu'on} obtient par conséquent
\[
\unsure{v_{\,\xi}|v_{1}|\psi}(s')=u(s')_{\mathbf C}\circ\bigl(i(s')\mid\mu_2\bigr):\mu_{2\,\mathbf C}\longrightarrow G_{\mathbf C}^{\xi(s')}
\]
\deleted{D'autre part,} Ceci posé, il faut que $\pi^{\xi(s')}$ (polarisation de $G_{\mathbf R}^{\xi(s')}$ déduite de $\pi$ sur $G$ par transport) soit celle définie par les $\mu_{2\,\mathbf C}$. Par conséquent, pour $s'$ \unsure{variable} de façon que $\xi(s')$ \unsure{reste|soit} constant, les polarisations \inserted{de $G_{\mathbf R}^{\xi(s')}$} associées aux différents $v(s')$ doivent être identiques -- ce qui donne une condition non tautologique sur la distribution des \deleted{\unsure{$v$}} $u(s')$. [\unsure{Noter|Notons} cependant que l'ens des \deleted{\unsure{classes de}} polarisations d'un \unsure{grpe|gpe} réel $H$, défini en termes d'hom $\mu_{2\,\mathbf C}\to H_{\mathbf C}$, \margin{? à vérifier} est \uline{fini}, [si $G$ est commutatif en tous cas !], donc la restriction n'est pas tellement draconienne.

Enfin, conformément aux \inserted{conjectures} de Tate magnifiées, la connaissance des $u(s')$ \unsure{permet} de retrouver la structure d'effectivité dans la catégorie
\margin{184}
